\section{Growing order and outer-envelope accuracy}
\label{sec:growing}
We now specialize to \(1\le T\le2\) and
\(V=8\sqrt2L\). The fixed-order theorem cannot be used by
substituting \(p=p(X)\) into an unspecified \(C_p\).
We instead retain the pre-absorption error estimates and bound
their constants simultaneously.

\subsection{Factorial bounds for the actual finite constants}
Put
\begin{align}
 E_*&=\max\{1,\sqrt2\e^{R+Y^2/4+4},
                         4\sqrt2\e^{Y^2/4}\},\\
 S_*&=\max\{1,W+4,16\sqrt2\},\qquad C_*=\max(1,A,D).
\end{align}
The half-integer gamma recurrence gives
\(\Gamma((m+1)/2)/\sqrt\pi\le m!\).
The finite moment sums consequently imply
\[
 M_m,\mathsf D_m\le E_*m!S_*^m,\qquad
 M_m^t,\mathsf D_m^t\le E_*(m+2)!S_*^{m+2}.
\]
Grouping the sums in \eqref{eq:majorant-polynomials} gives
\begin{subequations}\label{eq:factorial-majorants}
\begin{align}
 \mathsf B_m[M]&\le E_*(8C_*S_*^2)^m(2m)!,\\
 \mathsf B_m[M^t]&\le E_*S_*^2(8C_*S_*^2)^m(2m+2)!,\\
 \max\{\mathsf H_p[\mathsf D],\mathsf H_p[\mathsf D^t]\}
 &\le E_*S_*^2(p+1)(2C_*S_*^2)^p(2p+2)!.
\end{align}
\end{subequations}
For example, in the first line bound
\((m+k)!\le(2m)!\), \(S_*^{m+k}\le S_*^{2m}\),
and \(\sum_{k=1}^m\binom{m-1}{k-1}/k!\le2^{m-1}\).
Also
\begin{equation}
 C_{2p}\le C_RQ_R^p p!,\quad
 C_R=C_\Gamma\frac{2^{7/4}}7\sqrt\pi\,1.1(8\pi)^{3/4},
 \quad Q_R=1.1^2(8\pi).
 \label{eq:rs-factorial}
\end{equation}
Let \(Q_*=\max(1,8C_*S_*^2,2C_*S_*^2,Q_R)\),
\(T_*=\max(1,E_*S_*^2,C_R)\). The common explicit majorant
\[
 \mathcal J(p)=T_*(p+2)Q_*^{p+1}(2p+4)!
\]
bounds each cost in \eqref{eq:factorial-majorants} and \(C_{2p}\).
In particular
\begin{equation}
 \log\mathcal J(p)=O((p+1)\log(p+2)).
 \label{eq:log-growth}
\end{equation}
This bound precedes any choice \(p=O(L)\).
It follows directly from \(n!\le n^n\), with constants fixed
by the chart. The original finite sums remain the operative
constants when checking a particular height.

\subsection{Keep the tails before normalizing}
Write
\[
 Q_p=48\{C_{2p}M_0+C_{\mathrm{pop}}\e\,\mathsf B_{p+1}[M]\}
       +C_{2p}M_0^t+C_{\mathrm{pop}}\e\,\mathsf B_{p+1}[M^t].
\]
The proof of \cref{thm:fixed}, before tail absorption, gives
\begin{align}
 \lVert h-F_X^{[p,2p]}\rVert_{\gamma;X;1,2}^{(3,1)}
 \le C_\chi\bigl\{&
 Q_pX^{-(p+3/4)}
 +49(A_{\mathrm{mod}}X^{P_0-7/4}
                   +A_{\mathrm{art},p}X^{1/4})\e^{-L^2}\nonumber\\
 &+49A_{\mathrm{far}}X^{-7/4}
                       \e^{\pi X/8-X^2/256}\bigr\}.
 \label{eq:unabsorbed}
\end{align}
This estimate has common structural onsets for \(1\le p\le cL\),
each fixed \(c>0\): replace the retained-order condition by
\(2cL\le(X-V)/(4\pi)\). All other structural gates are
independent of \(p\).

The old fixed-order absorption condition
\[
 L^2\ge(p+1)\log X+\Lp A_{\mathrm{art},p}
\]
fails already at \(p=\lfloor L\rfloor\), because
\((p+1)\log X>L^2\).
That failure concerns an unnecessarily strong intermediate bound.
The unabsorbed estimate \eqref{eq:unabsorbed} remains available.

\subsection{A stronger Gaussian tail with visible costs}
Fix \(0<\theta<1\), put \(\delta=(1-\theta)/2\) and \(t_R=1+R/4\).
For \(\abs q\ge V\),
\[
 -\frac{(q-\Rea d)^2}{8}+\abs q
 \le-\frac{q^2}{8}+t_R\abs q
 \le-\frac{\theta q^2}{8}-\frac{\delta q^2}{8}
                         +\frac{2t_R^2}{\delta}.
\]
Define
\begin{subequations}\label{eq:theta-moments}
\begin{align}
 \mathsf D_m^\theta
 &=\frac{\e^{Y^2/4+2t_R^2/\delta}}{\sqrt{4\pi}}\,
          2^m(8/\delta)^{(m+1)/2}\Gamma((m+1)/2),\\
 (\mathsf D_m^\theta)^t
 &=\mathsf D_m^\theta/2+\mathsf D_{m+2}^\theta/4.
\end{align}
\end{subequations}
Extracting \(\e^{-\theta V^2/8}=\e^{-16\theta L^2}\)
and integrating the remaining Gaussian proves the refined tail
bound. Define \(A_{\mathrm{mod}}^\theta,A_{\mathrm{art},p}^\theta\)
by \eqref{eq:tail-constants} with these moments.
Then \eqref{eq:unabsorbed} holds with these constants and
\(\e^{-16\theta L^2}\) replacing \(\e^{-L^2}\).

The cost of the sharper tail is explicit. Use
\[
 E_\theta=\max\{1,\sqrt2\e^{R+Y^2/4+4},
        \e^{Y^2/4+2t_R^2/\delta}\sqrt{2/\delta}\},
 \quad
 S_\theta=\max\{1,W+4,4\sqrt{2/\delta}\}
\]
in \eqref{eq:factorial-majorants}.
Thus \(\log^+A_{\mathrm{art},p}^\theta
=O_\theta((p+1)\log(p+2))\).
The parameter \(\theta\) is fixed before the height tends to
infinity; the constants are not uniform as \(\theta\to1\).

\begin{theorem}[Controlled growing order]\label{thm:growing}
Fix \(r,b>0\), \(c>0\), and \(a\) with \(1/8<a\le c\).
For \(p=\lfloor aL\rfloor\), \(K=2p\), and every
\[
 0<\eta<\min(a-1/8,127/8),
\]
the full frozen model on \(1\le T\le2\) satisfies, eventually,
\begin{equation}
 \normo{h-F_X^{[p,2p]}}\le\e^{-\eta L^2}.
 \label{eq:growing-result}
\end{equation}
In particular \(p=\lfloor L/2\rfloor\), \(K=2p\) gives
\(\normo{h-F_X^{[p,2p]}}\le\e^{-L^2/4}\).
The estimate retains both arms, every \(n\le N\), and the bounded
completion integral.
\end{theorem}
\begin{proof}
Choose a fixed \(\theta<1\) such that
\(\eta<16\theta-1/8\).
By \eqref{eq:log-growth}, the logarithm of the source/local
constant is \(O_a(L\log L)\).
After \eqref{eq:outer-conversion}, its exponent is at most
\[
 -(p+3/4)\log X+L^2/8+O_a(L\log L)
 =-(a-1/8)L^2+O_a(L\log L).
\]
The moderate and artificial tails have exponents at most
\[
 -(16\theta-1/8)L^2+O_{a,\theta}(L\log L).
\]
The remote true tail is smaller still, with a negative quadratic
term in \(X\).
The strict inequalities for \(\eta\) make these bounds summably
smaller than the target after a common onset.
The finite four-part allocation is given in \cref{app:growing-gates}.
Taking \(\theta=1/2\), \(a=1/2\), \(\eta=1/4\) proves the
stated concrete choice.
\end{proof}
The range is sufficient, not asserted sharp. A cap \(p\le cL\)
alone would include \(p=1\) and would not imply outer decay.
The bound is relative to the positive envelope, including near
zeros of the signed field. Selection of \(p(X)\) does not permit
setting \(X=\Rea z\) during differentiation or differentiating
the floor. Each chart remains frozen.

The earlier ansatz investigation already supplied a separate
growing-degree outer approximation \cite{ans05}.
The present result establishes the corresponding accuracy for
the fluctuation grading, with simultaneous factorial costs and
the same canonical completion. It does not claim that the first
outer approximation has only now been obtained.
